\documentclass[12pt,a4paper]{article}

% Core Packages
\usepackage[utf8]{inputenc}
\usepackage[T1]{fontenc}
\usepackage{geometry}
\usepackage{graphicx}
\usepackage{hyperref}
\usepackage{titlesec}
\usepackage{enumitem}
\usepackage{booktabs}
\usepackage{array}
\usepackage{xcolor}
\usepackage{fancyhdr}
\usepackage{setspace}
\usepackage{parskip}
\usepackage{caption}
\usepackage{float}
\usepackage{amsmath}
\usepackage{amssymb}
\usepackage{algorithm}
\usepackage{algpseudocode}

% Page geometry
\geometry{margin=1in}

% Colors
\definecolor{indigoblue}{RGB}{0, 51, 102}
\definecolor{sectioncolor}{RGB}{0, 76, 153}

% Hyperref setup
\hypersetup{
    colorlinks=true,
    linkcolor=indigoblue,
    urlcolor=indigoblue,
    citecolor=indigoblue
}

% Section formatting
\titleformat{\section}{\Large\bfseries\color{sectioncolor}}{\thesection}{1em}{}
\titleformat{\subsection}{\large\bfseries\color{indigoblue}}{\thesubsection}{1em}{}

% Header/Footer
\pagestyle{fancy}
\fancyhf{}
\fancyhead[L]{\small COL761: Data Mining}
\fancyhead[R]{\small Assignment 1 - Part 3}
\fancyfoot[C]{\thepage}
\renewcommand{\headrulewidth}{0.4pt}

\onehalfspacing

\begin{document}

% ==============================================================================
% COVER PAGE
% ==============================================================================
\begin{titlepage}
    \centering
    \vspace*{1cm}
    {\Huge\bfseries\color{indigoblue} COL761 Data Mining}\\[0.5cm]
    {\LARGE Assignment 1}\\[0.3cm]
    {\Large\textbf{Part 3: Graph Indexing}}\\[2cm]
    
    {\LARGE\bfseries Team VindyCS5}\\[0.5cm]
    
    \begin{tabular}{ll}
        \toprule
        \textbf{Name} & \textbf{Entry Number} \\
        \midrule
        Abhiram Dharme & 2022CS51139 \\
        Raj Chandak & 2022CS51643 \\
        Purushottam Sharma & 2022CS51635 \\
        \bottomrule
    \end{tabular}\\[1cm]
    
    \rule{\textwidth}{0.5pt}\\[0.3cm]
    {\large Submitted: February 2, 2026}
    \vfill
\end{titlepage}

\tableofcontents
\newpage

%==============================================================================
% PROBLEM STATEMENT
%==============================================================================
\section{Problem Statement}
\textit{[10 (report) + 35 (competitive scoring) marks]}

Given a database of graphs $\mathcal{D} = \{g_1, g_2, \ldots, g_n\}$ and a query graph $q$, the goal is to efficiently identify a \textbf{candidate set} $C_q \subseteq \mathcal{D}$ containing all graphs that potentially contain $q$ as a subgraph.

The performance metric is:
\[
s_q = \frac{|R_q|}{|C_q|}
\]
where $|R_q|$ is the number of true positives (graphs actually containing $q$) and $|C_q|$ is the candidate set size. \textbf{Larger scores indicate better pruning.}

\subsection{Key Insight}

The necessary condition for subgraph isomorphism states: if $q \subseteq g_i$, then every substructure (fragment) present in $q$ must also be present in $g_i$. By identifying \textbf{discriminative subgraph fragments}, we can efficiently prune the candidate set.

%==============================================================================
% APPROACH
%==============================================================================
\section{Approach: gSpan-Based Frequent Subgraph Mining}

Our approach uses \textbf{gSpan} (graph-based Substructure Pattern Mining) to mine frequent subgraph patterns from the database, followed by optimized feature selection using density filtering, coverage-based scoring, and diversity maximization.

\subsection{Mining Pipeline}

\begin{enumerate}
    \item \textbf{Frequent Subgraph Mining}: Use gSpan to discover all frequent connected subgraphs
    \item \textbf{Support Tracking}: Track which database graphs contain each pattern (using \texttt{where} parameter)
    \item \textbf{Optimized Selection}: Select exactly 50 features using density filtering, coverage optimization, and anti-correlation scoring
\end{enumerate}

\subsection{Adaptive Support Mining}

We use an adaptive strategy to mine patterns at multiple support levels:

\begin{table}[H]
\centering
\begin{tabular}{ccc}
\toprule
\textbf{Support Level} & \textbf{Min Support (rel.)} & \textbf{Strategy} \\
\midrule
Level 1 & 30\% & Most selective patterns \\
Level 2 & 25\% & High-confidence patterns \\
Level 3 & 20\% & Moderate frequency \\
Level 4 & 15\% & Lower frequency \\
Level 5 & 10\% & Rare patterns \\
Fallback & 5\% & Emergency mining \\
\bottomrule
\end{tabular}
\caption{Adaptive Support Levels for Pattern Mining}
\end{table}

The algorithm starts at 30\% support and stops when at least 50 candidate patterns are found, enabling \textbf{early stopping} for efficiency.

\subsection{Pattern Constraints}

\begin{itemize}
    \item \textbf{Minimum vertices}: 2 (at least one edge)
    \item \textbf{Maximum vertices}: 8 (computational constraint)
    \item \textbf{Connected graphs only}: Ensures semantic coherence
\end{itemize}

%==============================================================================
% FEATURE SELECTION
%==============================================================================
\section{Optimized Feature Selection}

\subsection{Selection Framework}

We employ a \textbf{four-stage} optimization approach:

\subsubsection{Stage 1: Density Filtering}

We filter out features based on their coverage density to eliminate weak discriminators:

\textbf{High-Density Filter (>55\%):} Features present in more than 55\% of graphs provide weak filtering power and are removed:
\[
\text{if } \frac{\text{support}(f)}{|\mathcal{D}|} > 0.55 \implies \text{reject } f
\]

\textbf{Low-Density Filter (<10\%):} Very rare features appearing in fewer than 10\% of graphs lack sufficient coverage:
\[
\text{if } \frac{\text{support}(f)}{|\mathcal{D}|} < 0.10 \implies \text{reject } f
\]

\textbf{Rationale}: The optimal sweet spot is 10-55\% coverage, balancing discriminative power with query coverage.

\subsubsection{Stage 2: Coverage-Based Scoring}

For surviving patterns, we compute a coverage score that peaks at 25\% density:

\[
\text{coverage\_score}(f) = \max\left(0.1, 1.0 - 2.0 \times \left|\frac{\text{support}(f)}{|\mathcal{D}|} - 0.25\right|\right)
\]

\textbf{Intuition}: 25\% coverage represents an optimal balance—frequent enough to help with queries, rare enough to filter effectively.

\subsubsection{Stage 3: Size-Weighted Scoring}

Larger patterns impose stricter structural constraints. Medium-sized patterns (4-6 edges) receive exponential bonuses:

\[
w_{\text{size}}(f) = 
\begin{cases}
e^{0.55 \times |E|} & \text{if } |E| \geq 6 \\
e^{0.65 \times |E|} & \text{if } |E| = 5 \\
e^{0.70 \times |E|} & \text{if } |E| = 4 \\
e^{0.50 \times |E|} & \text{if } |E| = 3 \\
e^{0.25 \times |E|} & \text{if } |E| = 2
\end{cases}
\]

\textbf{Design}: We prefer 4-6 edge patterns which balance specificity (smaller candidate sets) with generality (applicable to many queries).

\subsubsection{Stage 4: Greedy Selection with Anti-Correlation}

After base scoring, we select features greedily to maximize diversity:

\begin{algorithm}[H]
\caption{Optimized Greedy Feature Selection}
\begin{algorithmic}[1]
\Require Scored patterns $P$, target count $k=50$
\Ensure Selected features $S$
\State $S \gets \emptyset$, $\text{covered} \gets \emptyset$, $\text{seen\_hashes} \gets \emptyset$
\For{$i = 1$ to $k$}
    \State $\text{best} \gets \text{null}$, $\text{best\_gain} \gets -1$
    \For{each pattern $p$ in $P$}
        \If{$\text{hash}(p) \in \text{seen\_hashes}$}
            \State \textbf{continue}
        \EndIf
        \State $\text{new\_cov} \gets |\text{where}(p) \setminus \text{covered}|$
        \State $\text{overlap\_ratio} \gets \frac{|\text{where}(p) \cap \text{covered}|}{|\text{where}(p)|}$
        \State $\text{anti\_corr} \gets 1.0 - 0.6 \times \text{overlap\_ratio}$
        \State $\text{coverage\_bonus} \gets 1.2$ if $0.10 \leq \text{coverage}(p) \leq 0.50$ else $1.0$
        \State $\text{gain} \gets (\text{new\_cov} + 1) \times \text{anti\_corr} \times |E(p)|^{0.5} \times \text{coverage\_bonus}$
        \If{$\text{gain} > \text{best\_gain}$}
            \State $\text{best} \gets p$, $\text{best\_gain} \gets \text{gain}$
        \EndIf
    \EndFor
    \State $S \gets S \cup \{\text{best}\}$
    \State $\text{seen\_hashes} \gets \text{seen\_hashes} \cup \{\text{hash}(\text{best})\}$
    \State $\text{covered} \gets \text{covered} \cup \text{where}(\text{best})$
\EndFor
\State \Return $S$
\end{algorithmic}
\end{algorithm}

\subsection{Key Optimization Parameters}

\begin{table}[H]
\centering
\begin{tabular}{lll}
\toprule
\textbf{Parameter} & \textbf{Value} & \textbf{Rationale} \\
\midrule
High-density filter & 55\% & Filters weak discriminators \\
Low-density filter & 10\% & Ensures sufficient coverage \\
Optimal coverage & 25\% & Balances filtering \& generality \\
Anti-correlation penalty & 0.6 & Moderate diversity enforcement \\
Coverage bonus range & 10-50\% & Rewards versatile features \\
Size exponent & 0.5 & Balanced pattern size preference \\
\bottomrule
\end{tabular}
\caption{Optimized Feature Selection Parameters}
\end{table}

\textbf{Performance}: This configuration achieves \textbf{0.5030 precision on Mutagenicity} and \textbf{0.3889 precision on NCI-H23}, representing significant improvements over baseline approaches.

%==============================================================================
% FEATURE EXTRACTION
%==============================================================================
\section{Feature Extraction via Subgraph Isomorphism}

\subsection{Convert Phase}

Once discriminative features are identified, we convert each graph (database or query) into a \textbf{binary feature vector}:

\[
v_g = (b_1, b_2, \ldots, b_k) \in \{0,1\}^k
\]

where $b_j = 1$ if and only if feature graph $f_j$ is subgraph-isomorphic to $g$.

\subsection{Implementation}

We use \textbf{NetworkX} for subgraph monomorphism testing:

\begin{verbatim}
from networkx.algorithms import isomorphism
matcher = isomorphism.GraphMatcher(
    graph, feature,
    node_match=lambda n1, n2: n1['label'] == n2['label'],
    edge_match=lambda e1, e2: e1['label'] == e2['label']
)
is_present = matcher.subgraph_is_monomorphic()
\end{verbatim}

This performs non-induced subgraph matching as required by the assignment.

%==============================================================================
% CANDIDATE GENERATION
%==============================================================================
\section{Candidate Set Generation}

\subsection{Necessary Condition}

For a database graph $g_i$ to contain query $q$ as a subgraph, every feature present in $q$ must also be present in $g_i$:

\[
g_i \in C_q \iff v_q \leq v_{g_i} \quad \text{(component-wise)}
\]

\subsection{Algorithm}

\begin{algorithm}[H]
\caption{Candidate Set Generation}
\begin{algorithmic}[1]
\Require Database features $V_D$, query features $v_q$
\Ensure Candidate set $C_q$
\State $C_q \gets \emptyset$
\For{each database graph $g_i$ with features $v_{g_i}$}
    \If{$v_q \leq v_{g_i}$ (all query features present)}
        \State Add $g_i$ to $C_q$
    \EndIf
\EndFor
\State \Return $C_q$
\end{algorithmic}
\end{algorithm}

\subsection{Optimization}

The component-wise comparison $v_q \leq v_{g_i}$ is implemented using NumPy's vectorized operations:

\begin{verbatim}
if np.all(query_vector <= db_vector):
    candidates.append(graph_id)
\end{verbatim}

This allows batch processing of all database graphs efficiently.

%==============================================================================
% IMPLEMENTATION DETAILS
%==============================================================================
\section{Implementation Details}

\subsection{Pipeline Overview}

\begin{enumerate}
    \item \textbf{identify.py}: Mine discriminative subgraphs using gSpan
    \begin{itemize}
        \item Parse database graphs into NetworkX format
        \item Convert to gSpan input format (integer labels)
        \item Run gSpan with \texttt{where} tracking at adaptive support levels (30\%, 25\%, 20\%, ...)
        \item Parse gSpan output to extract patterns with support counts
        \item Apply optimized feature selection (density filtering + coverage scoring + anti-correlation)
        \item Save exactly 50 selected features in graph format
    \end{itemize}
    
    \item \textbf{convert.py}: Convert graphs to feature vectors
    \begin{itemize}
        \item Load discriminativesubgraphs from \texttt{identify.py}
        \item For each database/query graph, test subgraph isomorphism
        \item Output binary NumPy matrix (graphs $\times$ 50 features)
    \end{itemize}
    
    \item \textbf{candidate\_gen.py}: Generate candidate sets
    \begin{itemize}
        \item Load database and query feature matrices
        \item For each query, find graphs where $v_q \leq v_{g_i}$
        \item Output candidates in 0-based indexing format
    \end{itemize}
    
    \item \textbf{eval.py}: Evaluate precision (optional)
    \begin{itemize}
        \item Compare candidate sets against ground truth
        \item Compute precision scores
    \end{itemize}
\end{enumerate}

\subsection{Libraries Used}

\begin{table}[H]
\centering
\begin{tabular}{ll}
\toprule
\textbf{Library} & \textbf{Purpose} \\
\midrule
gspan-mining & Frequent subgraph mining with DFS code \\
NetworkX & Graph representation, manipulation, and subgraph monomorphism \\
pandas & Data frame operations (patched for compatibility) \\
NumPy & Feature matrix operations \\
\bottomrule
\end{tabular}
\caption{External Libraries}
\end{table}

\subsection{Key Implementation Notes}

\subsubsection{Label Conversion}

gSpan requires \textbf{integer labels}. We maintain bidirectional mappings:
\begin{itemize}
    \item \texttt{node\_label\_map}: string $\rightarrow$ int
    \item \texttt{edge\_label\_map}: string $\rightarrow$ int
\end{itemize}

These are used to convert the database to gSpan format and restore original labels after mining.

\subsubsection{Canonical Hashing}

To avoid duplicate patterns, we use a canonical hash:

\begin{verbatim}
hash = (num_nodes, num_edges, 
        sorted(node_labels), 
        sorted(edge_triples))
\end{verbatim}

This ensures topologically identical patterns are deduplicated.

\subsubsection{0-Based Indexing}

All output uses \textbf{0-based indexing} for both query IDs and graph IDs, as required by the assignment specification.

%==============================================================================
% DESIGN DECISIONS
%==============================================================================
\section{Design Decisions and Rationale}

\subsection{Why gSpan for Pattern Mining?}

gSpan provides several advantages over manual pattern enumeration:

\begin{itemize}
    \item \textbf{Completeness}: Discovers all frequent connected subgraphs systematically
    \item \textbf{DFS Code}: Uses canonical depth-first search codes to avoid duplicate patterns
    \item \textbf{Efficiency}: Prunes the search space using anti-monotonicity (if pattern is infrequent, all superpatterns are too)
    \item \textbf{Where Tracking}: Optionally tracks which graphs contain each pattern (critical for correlation analysis)
\end{itemize}

Compared to manually mining specific pattern types (paths, cycles), gSpan can discover \textbf{arbitrary topologies} that may be more discriminative for specific datasets.

\subsection{Why Optimized Multi-Stage Feature Selection?}

Our selection combines four complementary stages:

\begin{table}[H]
\centering
\begin{tabular}{lp{8cm}}
\toprule
\textbf{Stage} & \textbf{Purpose} \\
\midrule
Density Filtering & Removes features that are too common (>55\%) or too rare (<10\%) \\
Coverage Scoring & Prioritizes features near 25\% density for optimal filtering \\
Size Weighting & Favors medium-sized patterns (4-6 edges) for balanced specificity \\
Anti-Correlation & Ensures diverse features using moderate overlap penalty (0.6) \\
\bottomrule
\end{tabular}
\caption{Optimized Multi-Stage Selection Rationale}
\end{table}

\textbf{Single-heuristic limitations}:
\begin{itemize}
    \item Pure frequency: May select many highly overlapping or too-common patterns
    \item Pure entropy: Ignores pattern size and doesn't filter extremes
    \item Pure size: Large patterns may be too rare; small patterns lack specificity
    \item Pure diversity: May select weak features just for novelty
\end{itemize}

\subsection{Why Adaptive Support Levels?}

Starting at high support (30\%) and progressively lowering has two benefits:

\begin{enumerate}
    \item \textbf{Efficiency}: gSpan terminates early if enough patterns are found at high support
    \item \textbf{Quality}: High-support patterns (within the 10-55\% sweet spot) are more likely to appear in test queries (better generalization)
\end{enumerate}

\subsection{Why 0-Based Indexing?}

The assignment specification requires 0-based indexing for graph and query IDs. This aligns with:
\begin{itemize}
    \item Python's 0-based indexing convention
    \item NumPy array indexing
    \item Modern programming practices
\end{itemize}

\subsection{Why NetworkX for All Graph Operations?}

\begin{itemize}
    \item \textbf{Compatibility}: Seamless integration with gSpan-mining library output
    \item \textbf{Flexibility}: Rich API for graph manipulation and analysis
    \item \textbf{Correctness}: Well-tested subgraph monomorphism implementation
    \item \textbf{Simplicity}: Single library reduces dependency complexity
\end{itemize}

NetworkX provides sufficient performance for our feature extraction needs while maintaining code clarity.

%==============================================================================
% ASSUMPTIONS
%==============================================================================
\section{Assumptions}

\begin{enumerate}
    \item \textbf{Graph Format}: Input graphs are in the specified format with lines starting with \texttt{v} for vertices and \texttt{e} for edges.
    \item \textbf{Duplicate Removal}: The database is pre-processed to remove duplicate graphs before calling \texttt{identify.py}, as specified in the assignment.
    \item \textbf{Feature Count}: We use exactly \textbf{50 features} selected through advanced multi-heuristic scoring.
    \item \textbf{Non-induced Isomorphism}: We use non-induced subgraph isomorphism (a feature is present if it appears as a subgraph, not necessarily an induced subgraph).
    \item \textbf{Binary Features}: All features are binary (0/1), indicating presence or absence of a substructure.
    \item \textbf{0-Based Indexing}: All query and graph IDs are 0-based, as required by the assignment specification.
    \item \textbf{Connected Patterns}: gSpan mines only connected subgraphs, ensuring semantic coherence of features.
\end{enumerate}

%==============================================================================
% CONCLUSION
%==============================================================================
\section{Conclusion}

Our gSpan-based graph indexing approach combines:
\begin{enumerate}
    \item \textbf{Systematic pattern mining}: gSpan discovers all frequent connected subgraphs with DFS-based canonicalization
    \item \textbf{Optimized multi-stage selection}: Density filtering (10-55\%) + coverage scoring (25\% target) + anti-correlation (0.6 penalty) + size balancing (0.5 exponent)
    \item \textbf{Exactly 50 features}: Compact index achieving 0.5030 (Mutagenicity) and 0.3889 (NCI-H23) precision
    \item \textbf{Efficient implementation}: NetworkX for mining and subgraph matching
\end{enumerate}

This approach generates compact candidate sets by:
\begin{itemize}
    \item Exploiting the necessary condition for subgraph monomorphism
    \item Filtering extreme densities (>55\% and <10\%) early to focus on sweet spot
    \item Prioritizing 25\% coverage for optimal balance
    \item Using moderate anti-correlation (0.6) to ensure diversity without over-penalizing useful overlap
    \item Balancing pattern size with 0.5 exponent to avoid bias toward large or small patterns
\end{itemize}

The adaptive support mining strategy (30\% $\rightarrow$ 25\% $\rightarrow$ 20\% ...) ensures we find high-quality patterns efficiently, while the four-stage optimized selection (density filtering, coverage scoring, size weighting, anti-correlation) maximizes both the discriminative power and diversity of the selected 50 features.

%==============================================================================
% REFERENCES
%==============================================================================
\newpage
\begin{thebibliography}{9}

\bibitem{yan2005gspan}
Yan, X., \& Han, J. (2002).
\textit{gSpan: Graph-based substructure pattern mining}.
IEEE ICDM, 721--724.

\bibitem{yan2004gindex}
Yan, X., Yu, P. S., \& Han, J. (2004).
\textit{Graph indexing: A frequent structure-based approach}.
ACM SIGMOD, 335--346.

\bibitem{cheng2007fg}
Cheng, J., Ke, Y., Ng, W., \& Lu, A. (2007).
\textit{Fg-index: towards verification-free query processing on graph databases}.
ACM SIGMOD, 857--872.

\bibitem{shang2008gsting}
Shang, H., Zhang, Y., Lin, X., \& Yu, J. X. (2008).
\textit{Taming verification hardness: An efficient algorithm for testing subgraph isomorphism}.
VLDB, 1--12.

\end{thebibliography}

\end{document}
